


\newacronym[
description={\todo}
]{rmsystem}{R/M-System}{Restriktions-/Modifikations-System}
\newglossaryentry{vorlage}{
name={Transposase},
description={\todo}
}
\newacronym{sam}{SAM}{S-Adenosyl-Methionin}
\newglossaryentry{replikase}{
name={Replikase},
description={\todo}
}
\newglossaryentry{Nitrocellulose-Filter}{
name={Nitrocellulose-Filter},
description={\todo}
}
\newglossaryentry{transkonjugant}{
name={Transkonjugant},
description={\todo}
}
\newglossaryentry{transferfunktion}{
name={Transferfunktion},
description={\todo}
}
\newglossaryentry{lbmedium}{
name={LB-Medium},
description={\todo}
}
\newglossaryentry{tymedium}{
name={TY-Medium},
description={\todo}
}
\newglossaryentry{typlatte}{
name={TY-Platte},
description={Siehe \gls{tymedium}}
}
\newacronym{mob}{mob}{Mobilisierungsregion}
\newglossaryentry{suicideplasmid}{
name={Suicide Plasmid},
description={Plasmide, welche in einem bestimmten Wirtsbakterium nicht repliziert werden können}
}
\newglossaryentry{transposase}{
name={Transposase},
description={desc}
}
\newglossaryentry{iselement}{
name={IS-Element},
description={desc}
}
\newacronym{ir}{IR}{inverse Sequenzwiederholung}
\newacronym{dr}{DR}{direkte Wiederholung}
\newglossaryentry{Isoform}{
name={Isoform},
description={Gen-Isoformen sind mRNAs, die vom gleichen Ort produziert werden, sich aber in ihren \glspl{tss}, \glspl{cds} und/oder \glspl{utr} unterscheiden, was die Genfunktion verändern kann}
}
\newacronym[
longplural={nichttranslatierten Regionen}
]{utr}{UTR}{nichttranslatierte Region}
\newacronym[
longplural={proteincodierenden DNA-Sequenzen}
]{cds}{CDS}{proteincodierende DNA-Sequenz}
\newacronym[
longplural={Transkriptionsstartstellen}
]{tss}{TSS}{Transkriptionsstartstelle}
\newacronym{rpkm}{RPKM}{Reads Per Kilobase of gene per Millions of total read}
\newacronym{repe}{repE}{Gen für Plasmid-Replikation}
\newacronym[
description={Künstlich geschaffenes DNA-Oligonukleotid aus rund 50bp, welches in ein Plasmid eingesetzt wurde und dessen Sequenz direkt hintereinander verschiedene Restriktionsschnittstellen für Restriktionsendonukleasen enthält. Bietet die Möglichkeit, mit Hilfe von Restriktionsenzymen Fremd-DNA in das Plasmid einzubauen}
]{mcs}{MCS}{Multiple Cloning Site}
\newglossaryentry{Operon}{
name={Operon},
description={Funktionseinheit der DNA von Prokaryoten und manchen Eukaryoten sowie der von Bakterien abgeleiteten Organellen wie den Plastiden. Besteht aus Promotor, \gls{Operator}(en) und mehreren (Struktur-)Genen, die für Proteine mit typischerweise verwandten Funktionen codieren. Je nach Operon können verschiedene regulatorische Proteine als Repressoren oder als Aktivatoren mit den Operatoren in Wechselwirkung treten und dadurch die Transkription der Gene im Operon an- oder abschalten}
}
\newglossaryentry{Operator}{
name={Operator},
description={Abschnitt auf dem \gls{Operon} in der Nähe oder innerhalb des Promotors, an den ein Regulatorprotein (Repressor oder Aktivator) binden kann und dadurch die Affinität des Promotors zur RNA-Polymerase verringert bzw. erhöht}
}
\newglossaryentry{sticky ends}{
name={sticky ends},
description={Enden, wenn ein Restriktionsenzym die DNA an versetzten Positionen auf den beiden Strängen schneidet, wodurch einsträngige Überhänge entstehen}
}
\newglossaryentry{blunt ends}{
name={blunt ends},
description={Enden, wenn ein Restriktionsenzym beide DNA-Stränge an genau gegenüberliegenden Positionen innerhalb der Erkennungssequenz schneidet}
}
\newglossaryentry{Trypsin}{
name={Trypsin},
description={Gemisch dreier Verdauungsenzyme, die im Dünndarm Proteine zersetzen und zu den Peptidasen zählen}
}
\newglossaryentry{Transkriptionseinheit}{
name={Transkriptionseinheit},
description={Gruppe von Genen, die unter der Kontrolle eines einzigen Promotors gemeinsam transkribiert werden}
}
\newglossaryentry{Phenol}{
name={Phenol},
description={, saures: Unterstützt die Entfernung von DNA (wird in die organische Phase abgeschieden), trägt zur Stabilisierung der Grenzschicht bei und verhindert Schaumbildung beim Mischen in RNA-Extraktionsverfahren}
}
\newglossaryentry{Ionenstärke}{
name={Ionenstärke},
description={(einer Lösung) Maß für die elektrische Feldstärke aufgrund gelöster Ionen}
}
\newglossaryentry{Corepressor}{
name={Corepressor},
description={Substanz, die die Expression von Genen hemmt.  Vermindert die Expression von Genen nicht durch direkte Interaktion mit dem \gls{Operator}, sondern indirekt durch Interaktion mit Proteinen, die sich ihrerseits an das Operatorgen binden. }
}
\newglossaryentry{Transcription terminator}{
name={Transkriptionsterminator},
description={Markiert das Ende eines Gens oder \gls{Operon}s}
}
\newglossaryentry{GC-Skew}{
name={GC-Skew},
description={Die Asymmetrie der Häufigkeit von Guanin (G) und Cytosin (C) in einem DNA-Strang relativ zueinander. In Bakterien kann der GC-Skew dazu verwendet werden, den Ursprung der DNA-Replikation und den Terminus zu identifizieren, da sich diese Regionen oft durch signifikante Veränderungen im GC-Skew-Profil auszeichnen}
}
\newglossaryentry{Glimmer}{
name={Glimmer},
description={Software system for finding genes in microbial DNA, especially the genomes of bacteria, archaea, and viruses}
}
\newglossaryentry{Evidenz}{
name={Evidenz},
description={Die verschiedenen Arten von Daten und Informationen, die zur Identifikation, Validierung und Charakterisierung von Genen verwendet werden}
}
\newglossaryentry{Inducer}{
name={Inducer},
description={Molekül, das die Transkription regulierter Gene steigert}
}
\newglossaryentry{16S-rRNA-Verwandschaft}{
name={16S-rRNA-Verwandschaft},
description={Kontext: Nutzung der 16S ribosomalen RNA (rRNA) als molekularen Marker zur Untersuchung der phylogenetischen Beziehungen und taxonomischen Klassifikation von Bakterien und Archaeen}
}
\newglossaryentry{Leseweite}{
name={Leseweite},
description={(Read Length) Die Länge eines DNA- oder RNA-Fragments, das während eines einzelnen Sequenzierungsvorgangs bestimmt wird}
}
\newglossaryentry{Trägerampholyt}{
name={Trägerampholyt},
plural={Trägerampholyte},
description={Mischung aus synthetischen amphoteren Molekülen, die in der Isoelectric Focusing (IEF) verwendet wird, einer Technik zur Trennung von Proteinen oder Peptiden basierend auf ihrem isoelektrischen Punkt}
}
\newglossaryentry{Streptavidin}{
name={Streptavidin},
description={Protein, das eine starke und spezifische Bindung zu Biotin (Vitamin B7) eingeht. Wird zur Bindung von Stoffen genutzt}
}
\newglossaryentry{Flow Cell}{
name={Flow Cell},
plural={Flow Cells},
description={Eine Flow Cell besteht aus einer kleinen, flachen Kammer aus Glas oder Kunststoff, die mit feinen Mikrokanälen durchzogen ist. Diese Kanäle ermöglichen den kontrollierten Fluss von Flüssigkeiten und Reagenzien während des Sequenzierungsprozesses}
}
\newacronym[
description={\textit{Posttranslationale Modifikation}. Veränderungen von Proteinen, die nach der Translation stattfinden. Die meisten werden durch den Organismus oder durch die Zellen selbst ausgelöst},
longplural={Posttranslationale Modifikationen}
]{ptm}{PTM}{Posttranslationale Modifikation}
\newacronym[
description={\textit{Enzyme comission number}. Numerisches Klassifikationssystem für Enzyme. Jede EC-Nummer besteht aus vier durch Punkte voneinander getrennten Zahlen. Genaugenommen wird nicht das Enzym selbst, sondern die Reaktion, die es katalysiert, kategorisiert. So kann es vorkommen, dass nicht miteinander verwandte Enzyme, welche die gleiche Reaktion katalysieren, dieselbe EC-Nummer zugewiesen bekommen}
]{ecnummer}{E.C Nummer}{Enzyme comission number}
\newacronym[
description={\textit{Expressed Sequence Tag}. Kurze DNA-Sequenzen, die von einem einzelsträngigen cDNA-Molekül stammen, das aus mRNA isoliert wurde. Sie repräsentieren Teile von Genen, die aktiv exprimiert werden}
]{est}{EST}{Expressed Sequence Tag}
\newacronym[
longplural={Didesoxyribonukleosid-Triphosphate},
description={\textit{Didesoxyribonukleosid-Triphosphat}. Artifizielle DNA-Nukleotide, die bei der DNA-Sequenzierung nach Sanger Verwendung finden. Sind wie \glspl{dntp} ausgebaut, allerdings ist die Ribose (Zucker) an Position 2’ und 3’ desoxidiert. Dadurch fehlt am 3'-Kohlenstoff-Atom die Hydroxygruppe, an der bei der Polymerisation das nächste Nukleotid angehängt wird}
]{ddntp}{ddNTP}{Didesoxyribonukleosid-Triphosphat}
\newglossaryentry{Scaffold}{
name={Scaffold},
description={Zusammensetzung von Contigs mithilfe von Paired-End-Bibliotheken}
}
\newglossaryentry{Polylinker}{
name={Polylinker},
description={Siehe \gls{mcs}}
}
\newglossaryentry{X-Gal}{
name={X-Gal},
description={Enzym. Kann von $\beta$-Galactosidase gespalten werden und daraus entsteht ein blauer Farbstoff. So
kann man $\beta$-Galactosidase nachweisen}
}


\newacronym{srna}{sRNA}{small RNA}
\newacronym{mirisc}{miRISC}{microRNA-induced silencing complex}
\newacronym{risc}{RISC}{RNA-induced silencing complex}
\newacronym{sirna}{siRNA}{short interfering RNA}
\newacronym{mirna}{miRNA}{microRNA}
\newacronym{pirna}{piRNA}{PIWI interacting RNA}
\newacronym{rbs}{RBS}{Ribosome Binding Site}
\newacronym{empcr}{emPCR}{Emulsion-PCR}
\newacronym{cogs}{COGc}{Clusters of Orthologous Groups}
\newacronym{sstdna}{sstDNA}{single stranded template DNA}
\newacronym{gfp}{GFP}{Grün fluoreszierendes Protein}
\newglossaryentry{sds}{
name={SDS-PAGE},
description={(sodium dodecyl sulfate polyacrylamide gel electrophoresis).  Analytische Methode der Biochemie zur Trennung von Stoffgemischen nach der Molekülmasse in einem elektrischen Feld}
}
\newacronym{hplc}{HPLC}{High Performance Liquid Chromatography}
\newacronym{nmr}{NMR}{Nuclear Magnetic Resonance}
\newacronym{lc}{LC}{Liquid Chromatography}
\newacronym{ce}{CE}{Kapillarelektrophorese}
\newacronym{gc}{GC}{Gaschromatographie}
\newacronym{ms}{MS}{Massenspektrometrie}
\newacronym{msms}{MSMS}{Tandem Massenspektrometrie}
\newglossaryentry{C-Wert}{
name={C-Wert},
description={Die Gesamtmenge an DNA in einem haploiden Genom}
} 
\newglossaryentry{Transposon}{
name={Transposon},
plural={Transposonen},
description={Mobile genetische Elemente und können sich von einem Ort im Genom zu einem anderen bewegen, entweder durch einen "Kopieren-und-Einfügen"-Mechanismus oder durch einen "Ausschneiden-und-Einfügen"-Mechanismus}
} 
\newglossaryentry{Lagging Strand Copy}{
name={Lagging Strand Copy},
description={Die diskontinuierlich synthetisierte Kopie eines DNA-Strangs}
} 
\newglossaryentry{Detergenzie}{
name={Detergenzie},
description={Chemische Substanzen, die verwendet werden, um Membranen zu lysieren, Proteine zu denaturieren und die Solubilisierung von lipophilen (fettlöslichen) Molekülen zu erleichtern}
} 
\newglossaryentry{Trichloressigsäure}{
name={Trichloressigsäure},
description={Starkes Denaturierungsmittel, das Proteine effizient präzipitiert. Dies ist besonders nützlich, um Proteine aus einer Probe zu entfernen und somit die Interferenzen in der nachfolgenden Analyse von Metaboliten zu minimieren}
}
\newglossaryentry{Assay}{
name={assay},
plural={assays},
description={Experimentelle Methode oder ein Verfahren, das verwendet wird, um die Anwesenheit, Menge oder Aktivität eines bestimmten biologischen Moleküls oder einer biologischen Funktion zu messen}
}
\newglossaryentry{Alkaloid}{
name={Alkaloid},
plural={Alkaloide},
description={Vielfältige Gruppe von natürlich vorkommenden organischen Verbindungen, die Stickstoff enthalten und meist basisch sind. Sie kommen in zahlreichen Pflanzenarten sowie in einigen Pilzen und Tieren vor}
}
\newglossaryentry{Substrat}{
name={Substrat},
description={Hier: Chemische Verbindung oder ein Molekül, das an ein Enzym bindet, um eine biologische Reaktion auszulösen oder zu beschleunigen}
}
\newglossaryentry{Endpunktextinktion}{
name={Endpunktextinktion},
description={Messung der Absorption oder \gls{Extinktion} einer Probe am Ende einer chemischen Reaktion in einem spektrophotometrischen Assay}
}
\newglossaryentry{Extinktion}{
name={Extinktion},
description={Die Verringerung der Intensität von Licht, das durch eine Lösung geleitet wird, die ein absorbierendes Material enthält. Maß für die Lichtabsorption durch die Lösung und wird oft in spektrophotometrischen Analysen verwendet, um die Konzentration von Substanzen zu bestimmen}
}
\newglossaryentry{Absorptionsspektrum}{
name={Absorptionsspektrum},
plural={Absorptionsspektren},
description={Wie viel Licht bei verschiedenen Wellenlängen durch eine Substanz absorbiert wird. Es ist ein charakteristisches Merkmal von Molekülen, Atomen oder Materialien und dient als wichtiger analytischer Fingerabdruck, um die chemische Zusammensetzung und Struktur einer Substanz zu identifizieren}
}
\newacronym[
description={\textit{Kyoto Encyclopedia of Genes and Genomes}. Sammlung von Datenbanken zu Genomen, biologischen Stoffwechselwegen, Krankheiten, Medikamenten und chemischen Substanzen}
]{kegg}{KEGG}{Kyoto Encyclopedia of Genes and Genomes}
\newacronym[
longplural={ Desoxyribonukleosid-Triphosphate},
description={\textit{ Desoxyribonukleosid-Triphosphat}. Nukleotid-Einheit aus Desoxyribose-Zucker, einer Phosphatgruppe und einer der vier stickstoffhaltigen Basen (Adenin, Thymin, Cytosin oder Guanin)}
]{dntp}{dNTP}{ Desoxyribonukleosid-Triphosphat}
\newacronym[
plural={RNasen},
description={Enzyme, die die hydrolytische Spaltung von Phosphodiesterbindungen in Ribonukleinsäure(RNA)-Ketten katalysieren. Diese Reaktion ist Teil der Prozessierung und des Abbaus von RNA}
]{RNase}{RNase}{Ribonuklease}

\newglossaryentry{Elutionspuffer}{
name={Elutionspuffer},
description={Lösung, um gebundene Moleküle von einem festen Träger (z.B. Mikrokügelchen) oder einer Matrix zu lösen und freizusetzen}
}
\newglossaryentry{Oligonukleotid}{
name={Oligonukleotid},
plural={Oligonukleotide},
description={Kurze, lineare Nukleinsäuresequenzen. Können DNA oder RNA sein.}
}
\newglossaryentry{Hexa-Nonanucleotid}{
name={Hexa-Nonanucleotid},
plural={Hexa-Nonanucleotide},
description={Der Begriff setzt sich aus Hexa (6), Nona (9) und Nukleotid zusammen, was bedeutet, dass die Oligonukleotide in dieser Kategorie typischerweise zwischen 6 und 9 Basenpaaren lang sind}
}
\newglossaryentry{Oligo-dT-Sequenz}{
name={Oligo-dT-Sequenz}, 
plural={Oligo-dT-Sequenzen}, 
description={Kurze DNA- oder RNA-Sequenzen, die aus wiederholten Thymin-Nukleotiden bestehen. Sie können als Primer für die cDNA Synthese, zum Binden der Poly-A-Schwänze von mRNA und in weiteren Hybridisierungsverfahren verwendet werden}
}
\newglossaryentry{Oligo-dT-Sonde}{
name={Oligo-dT-Sonde}, 
plural={Oligo-dT-Sonden}, 
description={\textit{Siehe \gls{Oligo-dT-Sequenz}.}}
}
\newacronym[
description={\textit{Guanidinisothiocyanatlösung}. Reagenz, das zur Zelllyse und zur Isolierung von RNA verwendet wird, indem es die Struktur von Biomolekülen denaturiert und gleichzeitig die RNA vor Abbau schützt}
]{gtc}{GTC}{Guanidinisothiocyanatlösung}
\newglossaryentry{RNA-Hydrolyse}{
name={RNA-Hydrolyse}, 
description={Der enzymatische oder chemische Prozess, bei dem RNA-Moleküle durch die Spaltung der Phosphodiesterbindungen zwischen den Nukleotiden abgebaut werden}
}
\newacronym[
longplural={3’-untranslatierten Regionen}
]{3utr}{3'UTR}{3’-untranslatierte Region}
\newglossaryentry{fplasmid}{
description={Plasmid, das Bakterien die Fähigkeit zur Konjugation (horizontaler Gentransfer) verleiht. Ermöglicht einen gerichteten Gentransfer vom Spender (dieser besitzt den F-Faktor, wird auch als F+ bezeichnet) zum Empfänger (F-)}, 
name={F-Plasmid},
plural={F-Plasmide}
}
\newglossaryentry{Gen-Silencing}{
description={ Beim Gen-Silencing erfolgt die Genregulation durch eine Hemmung der Übertragung (Transkription) einer genetischen Information von der DNA auf die mRNA (transkriptionelles Gen-Silencing) oder der nachfolgenden Übersetzung (Translation) der auf der mRNA gespeicherten Information in ein Protein (post-transkriptionelles Gen-Silencing)}, 
name={Gen-Silencing}
}
\newglossaryentry{RNA-Interferenz}{
description={Form des posttranskriptionelles Gen-Silencing. Ein natürlicher Mechanismus in den Zellen Eukaryoten, welcher der zielgerichteten Abschaltung von Genen dient. Sie ist ein Spezialfall des Gen-Silencing}, 
name={RNA-Interferenz}
}
\newglossaryentry{RNA-Virus}{
name={RNA-Virus},
description={Als RNA-Virus bezeichnet man Viren, deren Erbmaterial aus RNA besteht}
}
\newglossaryentry{Tag}{
description={Kurze Sequenzen von Nukleotiden, die spezifische Fragmente von Messenger-RNA (mRNA) repräsentieren}, 
name={Tag},
plural={Tags}
}
\newglossaryentry{photolithografisch Methode}{
description={Licht wird verwendet, um bestimmte chemische Reaktionen zu aktivieren, die zum Hinzufügen von Nukleotiden zur wachsenden DNA-Kette führen}, 
name={photolithografisch Methode},
plural={photolithografische Methoden}
}
\newglossaryentry{G-Strang-Überhang}{
name={G-Strang-Überhang},
description={Einzelsträngiger DNA-Überhang, der aus Guanin-Nukleotiden besteht und sich am Ende eines doppelsträngigen DNA-Moleküls befindet. Er tritt auf, wenn das 3'-Ende eines DNA-Strangs über das komplementäre Ende hinausragt, typischerweise am 3'-Ende eines doppelsträngigen DNA-Moleküls}
}
\newglossaryentry{tetracyclin}{
name={Tetracyclin},
description={Antibiotikum, das zur Gruppe der Tetracycline gehört}
}
\newglossaryentry{taq}{
name={Taq-Polymerase}, 
description={Thermostabile \gls{DNA-Polymerase} des Bakteriums Thermus aquaticus (Taq). Das Enzym wird hauptsächlich zur DNA-Vervielfältigung mit Hilfe der Polymerase-Kettenreaktion eingesetzt}
}
\newglossaryentry{Reverse Transkriptase}{
name={Reverse Transkriptase}, 
description={ Enzym, das als \gls{Polymerase} die Synthese von komplementärer DNA (cDNA) aus einer RNA-Vorlage katalysiert}
}
\newglossaryentry{DNA-Polymerase}{
name={DNA-Polymerase}, 
description={Enzyme, die als \gls{Polymerase} die Synthese von DNA aus Desoxyribonukleotiden katalysieren}
}
\newglossaryentry{Polymerase}{
name={Polymerase}, 
description={In allen Lebewesen vorkommende Enzyme, die die Polymerisation von Nukleotiden, den Grundbausteinen der Nukleinsäure, katalysieren}
}
\newglossaryentry{telomerase}{
description={ Enzym des Zellkerns, welches aus einem Protein- (TERT) und einem langen RNA-Anteil (TR) besteht und somit ein Ribonukleoprotein ist. Dieses Enzym stellt die Endstücke der Chromosomen, die sogenannten Telomere, wieder her}, 
name={Telomerase},
plural={Telomerasen}
}
\newglossaryentry{okazaki}{
description={kurze DNA-Stücke, die während der Replikation des lagging strands synthetisiert werden. Da die DNA-Polymerase nur in 5'-zu-3'-Richtung arbeiten kann, wird der lagging strand diskontinuierlich repliziert. Die Okazaki-Fragmente müssen durch RNA-Primer initiiert werden},
name={Okazaki-Fragment}
}
\newglossaryentry{vektor}{
description={Transportvehikel/ eine Transportmatrix zur Übertragung einer Fremd-Nukleinsäure (oft DNA) in eine lebende Empfängerzelle durch Transfektion oder Transformation}, 
name={Vektor},
plural={Vektoren}
}
\newglossaryentry{haarnadelstruktur}{
description={Sekundäre Stamm-Schleife-Strukturen (englisch stem-loop) mit doppelsträngig ausgebildetem Stamm und kurzer einzelsträngiger Schleife. Diese treten auf, wenn zwei Abschnitte des gleichen Moleküls – oft mit palindromischer Nukleotid-Sequenz – durch komplementäre Basenpaare eine doppelsträngige Region bilden und den ungepaarten Zwischenabschnitt als Schleife einschließen}, 
name={Haarnadelstruktur},
plural={Haarnadelstrukturen}
}
\newglossaryentry{telomer}{
description={Die aus repetitiver DNA und assoziierten Proteinen bestehenden Enden der linearen Chromosomen}, 
name={Telomer},
plural={Telomere}
}
\newglossaryentry{bChromosom}{
description={Bakterielle Chromosomen enthalten Housekeeping-Gene, die für das Überleben der Zelle unerlässlich sind},
name = {bakterielles Chromosom}
}
\newglossaryentry{replikon}{
name={Replikon},
description={DNA- oder RNA-Abschnitt, der ausgehend von einem Replikationsursprung an einem Stück repliziert wird},
plural={Replikons}
}
\newglossaryentry{plasmid}{
description={Plasmide sind extrachromosomale Elemente, die in den meisten Bakterien vorkommen. Im Allgemeinen ist die genetische Information entbehrlich. Plasmide können der Bakterienzelle Wachstumsvorteile verschaffen, z. B. katabolische Aktivitäten oder Antibiotikaresistenzen}, 
name={Plasmid},
plural={Plasmide}
}
\newglossaryentry{megaplasmid}{
description={Megaplasmide sind große bakterielle Replikons, die einen plasmidartigen Replikationsursprung enthalten.
Einige Megaplasmide scheinen Gene zu enthalten, die für das Überleben der Zelle wesentlich sind.
Daher sind diese Replikone Grenzfälle. Bei einigen Megaplasmiden ist es schwierig zu entscheiden, ob sie als Chromosomen oder Plasmide bezeichnet werden sollten}, 
name={Megaplasmid}
}
\newglossaryentry{contig}{
name={Contig},
description={Ein Satz überlappender DNA- oder Protein-Stücke (reads), die von derselben genetischen Quelle stammen}
}
\newacronym[
description={\textit{Structural maintenace chromosomes}. }
]{smc}{SMC}{Structural maintenace chromosomes}
\newacronym{cid}{CID}{Chromosomal Interaction Domain}
\newacronym{hdr}{HDR}{High-Density Chromosomal region}
\newacronym[
description={Der Replikationsursprung im bakteriellen Chromosom bwz. eukaryotischen Genom}
]{oric}{OriC}{Origin of Chromosomal Replication}
\newacronym{terc}{TerC}{Termination of Replication}
\newacronym[
description={Der Replikationsursprung im Plasmid}
]{oriv}{OriV}{Origin of Vegetative Replication}
\newacronym[
longplural={bacterial artificial chromosomes},
plural={BACs}
]{bac}{BAC}{bacterial artificial chromosome}
\newacronym[
longplural={yeast artificial chromosomes},
plural={YACs}
]{yac}{YAC}{yeast artificial chromosome}
\newacronym[
longplural={Long interspersed nuclear elements},
plural={LINEs}
]{line}{LINE}{Long interspersed nuclear element}
\newacronym{orf}{ORF}{Open Reading Frame}
\newacronym[
longplural={Short interspersed nuclear elements},
plural={SINEs}
]{sine}{SINE}{Short interspersed nuclear element}
\newacronym{rflp}{RFLP}{Restriktionsfragmentlängen-Polymorphismus}
\newacronym{sslp}{SSLP}{Simple sequence length polymorphism}
\newacronym{snp}{SNP}{Single nucleotide polymorphism}
\newacronym{str}{STR}{simple tandem repeat}
\newacronym{vntr}{VNTR}{variable number of tandem repeat}
\newacronym{sybr}{SG}{SYBR® Green (SG)}
\newglossaryentry{pairedend}{
name={Paired-End-Read},
description={Ein Paar von sequenzierten Enden von kürzeren DNA-Fragmenten. Nützlich für die Assemblierung von Contigs und die Analyse von strukturellen Varianten in relativ kleinen Bereichen}
}
\newglossaryentry{matepair}{
name={Mate-Pair-Read},
description={Ein Paar von sequenzierten Enden von großen DNA-Fragmenten. So kann man weit entfernte Regionen in Kontext setzen -> Relevant für Scaffolding}
}
\newglossaryentry{Konsensussequenz}{
name={Konsensussequenz},
description={Sequenz von Nukleotiden oder Aminosäuren, welche in der Summe am wenigsten von einer gegebenen Menge von entsprechenden Mustersequenzen abweicht. Die genaue Beschaffenheit dieser Sequenz kann hierbei je nach Wahl des Abstandsmaßes, wie etwa Hamming- oder Levenshtein-Distanz, variieren}
}
\newacronym{olc}{OLC}{Overlap-layout-consensus}
\newglossaryentry{Haplotyp}{
name={Haplotyp},
plural={Haplotypen},
description={Die Kombination der Allele mehrerer gekoppelter Gene eines einzelnen Chromosoms bzw. Variante der Gesamtheit der Gene/der Nukleotidsequenz eines Chromosoms}
}
